%==============================================================================
\section{Phân cụm phân cấp (Hierarchical Clustering)}
\label{sec:hierarchical}
%==============================================================================

\begin{dinhnghia}[title={Phân cụm phân cấp}]
Phân cụm phân cấp là thuật toán học \textbf{không giám sát} gom các điểm \textbf{không nhãn} có đặc tính tương tự. Cấu trúc cụm được biểu diễn dạng \textbf{cây} hoặc \textbf{dendrogram}.
\end{dinhnghia}

\begin{phuongphap}[title={Hai họ thuật toán}]
\begin{itemize}
  \item \textbf{Agglomerative} (gộp từ dưới lên): mỗi điểm ban đầu là một cụm, lần lượt gộp cặp cụm gần nhất.
  \item \textbf{Divisive} (chia từ trên xuống): ban đầu một cụm lớn, lần lượt chia thành cụm nhỏ hơn.
\end{itemize}
\end{phuongphap}

\subsection{Các bước phân cụm phân cấp gộp (Agglomerative)}

\begin{phuongphap}[title={Quy trình}]
\begin{itemize}
  \item \textbf{Bước 1}: Coi mỗi điểm là một cụm --- ban đầu có $K$ cụm bằng số điểm.
  \item \textbf{Bước 2}: Gộp hai điểm \textbf{gần nhất} thành một cụm $\Rightarrow$ còn $K-1$ cụm.
  \item \textbf{Bước 3}: Tiếp tục gộp hai cụm gần nhất $\Rightarrow$ $K-2$ cụm.
  \item \textbf{Bước 4}: Lặp đến khi còn một cụm duy nhất (mọi điểm đã được gộp).
  \item \textbf{Bước 5}: Dùng \textbf{dendrogram} để \textbf{cắt} cây thành số cụm phù hợp bài toán.
\end{itemize}
\end{phuongphap}

\subsection{Vai trò dendrogram}

\begin{ynghia}[title={Dendrogram sau khi gộp hết}]
Khi đã có một cụm lớn duy nhất, dendrogram cho phép \textbf{tách} lại thành nhiều cụm con có quan hệ gần gũi, tùy ngưỡng cắt hoặc số cụm mong muốn.
\end{ynghia}

\subsubsection{Ví dụ --- Dendrogram}

\begin{vidu}[title={Import và 10 điểm 2D}]
\begin{lstlisting}
import matplotlib.pyplot as plt
import numpy as np

X = np.array([[7, 8], [12, 20], [17, 19], [26, 15], [32, 37],
              [87, 75], [73, 85], [62, 80], [73, 60], [87, 96]])
labels = range(1, 11)
plt.figure(figsize=(10, 7))
plt.subplots_adjust(bottom=0.1)
plt.scatter(X[:, 0], X[:, 1], label='True Position')
for label, x, y in zip(labels, X[:, 0], X[:, 1]):
    plt.annotate(label, xy=(x, y), xytext=(-3, 3),
                 textcoords='offset points', ha='right', va='bottom')
plt.show()
\end{lstlisting}
\end{vidu}

\begin{ynghia}[title={Hai nhóm trực quan}]
Trên đồ thị dễ thấy hai nhóm điểm; dữ liệu thực tế có thể có hàng nghìn cụm tiềm ẩn.

\tsimg{06_hierarchical/img01.jpg}
\end{ynghia}

\begin{vidu}[title={Vẽ dendrogram với SciPy}]
\begin{lstlisting}
from scipy.cluster.hierarchy import dendrogram, linkage

linked = linkage(X, 'single')
labelList = range(1, 11)
plt.figure(figsize=(10, 7))
dendrogram(linked, orientation='top', labels=labelList,
           distance_sort='descending', show_leaf_counts=True)
plt.show()
\end{lstlisting}
\end{vidu}

\begin{ynghia}[title={Dendrogram linkage single}]
\tsimg{06_hierarchical/img02.jpg}
\end{ynghia}

\begin{phuongphap}[title={Chọn số cụm bằng khoảng cách dọc lớn nhất}]
Chọn \textbf{đoạn dọc dài nhất} trên dendrogram, kẻ đường ngang cắt qua: số lần cắt cho biết số cụm. Ví dụ đường ngang cắt hai nhánh $\Rightarrow$ \textbf{2 cụm}.

\tsimg{06_hierarchical/img03.jpg}
\end{phuongphap}

\begin{vidu}[title={AgglomerativeClustering (sklearn)}]
\begin{lstlisting}
from sklearn.cluster import AgglomerativeClustering

cluster = AgglomerativeClustering(n_clusters=2, linkage='ward')
cluster.fit_predict(X)

plt.scatter(X[:, 0], X[:, 1], c=cluster.labels_, cmap='rainbow')
\end{lstlisting}
\end{vidu}

\begin{ynghia}[title={Hai cụm sau gán nhãn}]
\tsimg{06_hierarchical/img04.jpg}
\end{ynghia}

\begin{luuy}[title={Minh họa bổ sung}]
Các hình \texttt{img05}, \texttt{img06} minh họa thêm các bước gộp cụm và kết quả phân tách trên ví dụ tương tự.

\tsimg{06_hierarchical/img05.jpg}

\tsimg{06_hierarchical/img06.jpg}
\end{luuy}

\subsubsection{Ví dụ --- Cụm trên dữ liệu y khoa}

\begin{dongco}[title={Dataset tiểu đường}]
Ví dụ minh họa dùng \texttt{sklearn.datasets.load\_diabetes} (nhiều feature, bối cảnh y khoa). Có thể thay bằng \texttt{fetch\_openml(name='diabetes', version=1)} nếu cần bộ Pima 8 cột.
\end{dongco}

\begin{luuy}[title={Đường dẫn CSV cục bộ}]
Nếu có file Pima Indians Diabetes cục bộ, đặt đường dẫn hợp lệ và đọc bằng \texttt{pandas.read\_csv}; không bắt buộc khi đã dùng \texttt{load\_diabetes}.
\end{luuy}

\begin{vidu}[title={Nạp dữ liệu và chọn hai feature}]
\begin{lstlisting}
import matplotlib.pyplot as plt
import numpy as np
from sklearn.datasets import load_diabetes
import scipy.cluster.hierarchy as shc
from sklearn.cluster import AgglomerativeClustering

data = load_diabetes()
X_full = data.data
patient_data = X_full[:, 3:5]   # hai cột feature (tương tự chọn cột 3:5)
\end{lstlisting}
\end{vidu}

\begin{vidu}[title={Dendrogram trên toàn bộ feature}]
\begin{lstlisting}
plt.figure(figsize=(10, 7))
plt.title("Patient Dendograms")
dend = shc.dendrogram(shc.linkage(X_full, method='ward'))
\end{lstlisting}
\end{vidu}

\begin{vidu}[title={Phân 4 cụm trên hai chiều đã chọn}]
\begin{lstlisting}
cluster = AgglomerativeClustering(n_clusters=4, linkage='ward')
cluster.fit_predict(patient_data)

plt.figure(figsize=(7.2, 5.5))
plt.scatter(patient_data[:, 0], patient_data[:, 1],
            c=cluster.labels_, cmap='rainbow')
plt.show()
\end{lstlisting}
\end{vidu}

\begin{ynghia}[title={Hai đồ thị đầu ra}]
Chạy mã trên sinh dendrogram (toàn feature) và scatter hai chiều tô màu theo 4 cụm --- tương ứng hai hình minh họa trong bài (dendrogram bệnh nhân và phân cụm 2D).
\end{ynghia}

\begin{tinhchat}[title={Linkage}]
\texttt{single}, \texttt{complete}, \texttt{average}, \texttt{ward} ảnh hưởng cách đo khoảng cách giữa cụm; \texttt{ward} thường cho cụm compact khi feature đồng thang.
\end{tinhchat}

\begin{luuy}[title={Chi phí}]
Phân cụm phân cấp tốn $O(n^2)$--$O(n^3)$ bộ nhớ/thời gian; dataset rất lớn nên cân nhắc K-means hoặc BIRCH.
\end{luuy}
